% Section D: attitude estimation and fusion (LP14-LP17).
\section{Attitude estimation}

\begin{frame}{Two imperfect sensors, one good estimate \hfill \lp{14}}
  \begin{columns}[T]
    \begin{column}{0.48\textwidth}
      \footnotesize
      \begin{tabular}{@{}lll@{}}
        \toprule
        & Gyro $\int\omega$ & Accel tilt \\
        \midrule
        Fast motion & \textcolor{MediumGreen!60!black}{good} & noisy, spoofed \\
        Slow / DC & drifts & \textcolor{MediumGreen!60!black}{absolute} \\
        \bottomrule
      \end{tabular}
      \vspace{0.4em}

      \small Complementary filter:
      \[ \hat\theta_k = \alpha\,(\hat\theta_{k-1}+\omega_k\Delta t) + (1-\alpha)\,\theta^{\mathrm{acc}}_k \]
      \[ \hat\Theta = \tfrac{s\tau}{1+s\tau}\,\Theta_{\mathrm{gyro}} + \tfrac{1}{1+s\tau}\,\Theta_{\mathrm{acc}} \]
      \[ \tau=\tfrac{\alpha\Delta t}{1-\alpha},\qquad f_c = \tfrac{1}{2\pi\tau} \]
      \footnotesize Choose $\tau$ from the physics, then derive $\alpha$ from $\Delta t$.
    \end{column}
    \begin{column}{0.5\textwidth}
      \includegraphics[width=\linewidth]{cf_bode.pdf}
    \end{column}
  \end{columns}
  \tryit{complementary}{balanced\_sine, trust\_gyro, trust\_accel, step\_pitch}
\end{frame}

\begin{frame}{Kalman filter: estimate the bias too \hfill \lp{16}}
  \begin{columns}[T]
    \begin{column}{0.45\textwidth}
      \small
      State $\vect x=[\theta,\ b]^\top$, gyro as the control input:
      \[
        \vect x_k = \begin{bmatrix}1 & -\Delta t\\ 0 & 1\end{bmatrix}\vect x_{k-1} + \begin{bmatrix}\Delta t\\0\end{bmatrix}\tilde\omega_k + \vect w_k
      \]
      Accel tilt measurement: $z_k = \theta^{\mathrm{acc}}_k = [1\ 0]\,\vect x_k + v_k$
      \begin{itemize}
        \item $Q$: gyro noise and bias random walk. $R$: accel tilt noise.
        \item The covariance $P$ says how confident the filter is.
        \item Steady state: $K_\theta\to$ constant, so it is \emph{a complementary filter} with $\alpha_{\mathrm{eq}}=1-K_\theta$, \emph{plus} a bias estimate.
      \end{itemize}
    \end{column}
    \begin{column}{0.53\textwidth}
      \includegraphics[width=\linewidth]{cf_kf_compare.pdf}
    \end{column}
  \end{columns}
  \tryit{complementary}{kalman\_bias, kalman\_tuning}
\end{frame}

\begin{frame}{3-D: the Mahony nonlinear complementary filter \hfill \lp{15}}
  \begin{columns}[T]
    \begin{column}{0.5\textwidth}
      \small
      \[
        \dot{\vect q} = \tfrac12\,\vect q\qmul\big[0,\ \tilde{\boldsymbol\omega}-\hat{\vect b}+K_p\,\vect e\big],\qquad
        \dot{\hat{\vect b}} = -K_i\,\vect e
      \]
      \[
        \vect e = \hat{\vect a}\times\hat{\mat R}^\top\!\begin{bmatrix}0\\0\\-1\end{bmatrix}
        \;+\; \hat{\vect m}\times\hat{\mat R}^\top\vect b_{\mathrm{ref}}
      \]
      \begin{itemize}
        \item The cross product gives the rotation that aligns predicted and measured directions.
        \item $K_p$ plays the role of $1/\tau$. $K_i$ is integral action, which estimates the gyro bias.
        \item Runs at kHz on a microcontroller: the PX4/ArduPilot fallback, and many drones.
      \end{itemize}
    \end{column}
    \begin{column}{0.48\textwidth}
      \begin{block}{Madgwick, EKF, ...}
        \small Madgwick uses a gradient step for the same correction. An EKF/ESKF adds covariance and more states (velocity, position, accel bias). The structure is the same: \emph{predict with the gyro, correct with vector observations}.
      \end{block}
    \end{column}
  \end{columns}
  \tryit{complementary}{mahony\_with\_mag}
\end{frame}

\begin{frame}{Observability: gravity cannot see yaw \hfill \lp{16}}
  \begin{columns}[T]
    \begin{column}{0.58\textwidth}
      \includegraphics[width=\linewidth]{mahony_observability.pdf}
    \end{column}
    \begin{column}{0.4\textwidth}
      \small
      \begin{itemize}
        \item Gravity is invariant to rotation about the vertical. It corrects roll and pitch, and the $x,y$ gyro biases.
        \item Without a heading reference, yaw error $\approx \int b_z\,dt$: it grows forever, and $\hat b_z$ does not converge.
        \item Heading references: magnetometer, GNSS velocity, visual landmarks, loop closure.
        \item Ask of every filter state: \emph{does any measurement see it?}
      \end{itemize}
    \end{column}
  \end{columns}
  \tryit{complementary}{mahony\_no\_mag, mahony\_with\_mag}
\end{frame}

\begin{frame}{When the gravity assumption fails \hfill \lp{17}}
  \begin{columns}[T]
    \begin{column}{0.5\textwidth}
      \includegraphics[width=\linewidth]{surge_gating.pdf}
    \end{column}
    \begin{column}{0.48\textwidth}
      \small
      \begin{itemize}
        \item Every accel correction assumes $\vect f\approx-\vect g$. Acceleration during manoeuvres violates it.
        \item \textbf{Gating}: skip updates when $\big|\|\vect f\|-g\big|>\epsilon$. Cheap, but blind when tilt cancels the norm change.
        \item \textbf{Adaptive $R$}: inflate the measurement noise with $\big|\|\vect f\|-g\big|$ or the gyro rate.
        \item \textbf{Model it}: estimate velocity (INS + aiding) and subtract the true acceleration. This is what full navigation filters do.
      \end{itemize}
    \end{column}
  \end{columns}
  \tryit{complementary}{surge\_spoofs\_accel, surge\_gated, mahony\_surge}
\end{frame}

\begin{frame}{\checkq{} --- attitude estimation}
  \begin{enumerate}
    \item Your loop runs at 400 Hz with $\alpha=0.98$. What is $\tau$? What happens if you move the code to a 100 Hz loop unchanged?
    \item Which frequency band does each sensor contribute in a complementary filter, and why?
    \item How is a steady-state Kalman filter related to a complementary filter? What does it add?
    \item Why can't an accelerometer-plus-gyro filter estimate the vertical gyro bias on level ground?
    \item A car brakes hard for 3 s. Sketch the pitch error of a CF with and without gating.
    \item When would you prefer an EKF over Mahony on a robot?
  \end{enumerate}
\end{frame}
